\documentclass[11pt]{article}

% ---------------------------------------------------------------------------
% The F2Z forest-GKR note. Build: latexmk -pdf main.tex
% How the merged, lazy, bit-affine grand-product forest is PERFORMED by the
% prover in this repository (src/merged_forest.rs + the eq-factored driver).
% Companion notes: docs/verifier-note/ (the checks), docs/rlc-family-note/
% (the experimental claim families).
% ---------------------------------------------------------------------------

\usepackage[margin=1.05in]{geometry}
\usepackage{amsmath,amssymb,amsthm}
\usepackage{booktabs}
\usepackage{array}
\usepackage{enumitem}
\usepackage{microtype}
\usepackage{xcolor}
\usepackage[colorlinks=true,linkcolor=blue!50!black,citecolor=blue!50!black,urlcolor=blue!50!black]{hyperref}

% --- notation (aligned with docs/verifier-note/main.tex) --------------------
\newcommand{\F}{\mathbb{F}}
\newcommand{\Z}{\mathbb{Z}}
\newcommand{\K}{\mathbb{K}}            % the binding field GF(2^128)
\newcommand{\Kx}{\K^{\times}}
\newcommand{\Ftwo}{\F_2}
\newcommand{\Fq}{\F_q}                 % the evaluation field (char != 2)
\newcommand{\bits}{\{0,1\}}
\newcommand{\INT}{\mathsf{INT}}
\newcommand{\mle}[1]{\widetilde{#1}}
\newcommand{\eq}{\mathsf{eq}}
\newcommand{\cw}{c_w}
\newcommand{\albert}[1]{\textcolor{teal}{Albert: {#1}}}
% \code: typewriter names; underscores allow a line break after them.
\newcommand{\code}[1]{{\def\_{\textunderscore\allowbreak}\texttt{#1}}}

% Typewriter-heavy paragraphs: give the line breaker slack.
\tolerance=1600
\setlength{\emergencystretch}{1.8em}
\hbadness=2500

\newtheorem{definition}{Definition}
\newtheorem{remark}{Remark}
\newtheorem{identity}{Identity}

\title{How F2Z performs the forest GKR:\\
  the merged, lazy, bit-affine grand-product forest}
\author{The \code{f2z-pcs} repository}
\date{2026-07-27}

\begin{document}
\maketitle

\begin{abstract}
\noindent
F2Z's integer-MLE opening rides on a grand-product argument: per column
$c$ of the committed bit matrix, the prover must certify
$\alpha^{u_c} = \prod_i \Lambda[c][i]$ for the \emph{bit-affine} leaves
$\Lambda[c][i] = 1 + M[c][i]\cdot\tau(i)$.  This is the prover's dominant
phase (81--96\,\% of prove time across $n = 16$--$32$), and this note
explains how it is actually performed, as implemented in
\code{src/merged\_forest.rs} and the eq-factored sumcheck driver it runs
on.  Three ideas carry the design.  (i)~\emph{Merging}: the $2^s$
product trees are one multilinear family with the tree index as extra
MLE variables, so one GKR run binds every root and each layer's payload
is $O(\ell + s)$ field elements instead of $O(2^s)$.  (ii)~\emph{The
two-phase eq-factored layer sumcheck}: phase~A binds the in-tree
variables with the trees as eq-weighted groups sharing one reduction
point; phase~B binds the tree-index variables over the per-tree finals,
which never enter the transcript.  The driver materialises no $\eq$
tables, accumulates round coefficients in unreduced $256$-bit form, and
gets the fourth Lagrange node for free from a characteristic-2
affine-flat identity.  (iii)~\emph{Laziness}: the leaf layer and the
bottom product levels are never built --- the bottom layers run
branchless case-LUT sumcheck rounds straight off the packed committed
bits against shared $\tau$-product tables, one intermediate level is
regenerated just in time (with its first round message fused into the
generation pass), and the stored chain tops out a quarter (default
schedule) or an eighth (\code{l8}) of the way up the tree, which is what
caps prover memory at ${\approx}\,2^n{\cdot}4$\,B resp.\
${\approx}\,2^n{\cdot}2$\,B.  Every lazy path is an exact
characteristic-2 rewriting of the eager forest: the transcripts are
byte-identical, and tests pin this in both schedules.  The note closes
with the exit-claim discharge into the ring-switch/Ligerito opening, the
batched and RLC-family variants of the forest, and the measured cost
tables.
\end{abstract}

\tableofcontents

% ===========================================================================
\section{The job of the forest}
\label{sec:job}

F2Z\footnote{Repository \code{f2z-pcs}, a standalone extraction of the
\code{f2-int} pipeline from \code{zinc-plus} (branch
\code{f2-int-unified-ligerito}); the construction is due to Lev Soukhanov
(the char2-fieldswitch note, \S9).  \code{docs/DESIGN.md} holds the
distilled protocol; \code{docs/verifier-note/} specifies every verifier
check; \code{docs/rlc-family-note/} develops the experimental claim
families.  This note is the \emph{prover-side} companion: how the forest
is computed, scheduled, and kept small.} proves an \emph{integer}
multilinear-extension evaluation
$\smash{\mle{\INT(D)}}(r) = y \in \Fq$ over data $D$ whose \emph{bits}
are committed with a characteristic-2 commitment.  The instance shape is
$(t, s, W)$ (\code{IntEvalParams}): $D$ is a $2^t \times 2^s$ matrix of
$W$-bit cells, its bits arranged as the single-bit matrix
\[
  M[c][i] \in \bits, \qquad
  i = (b \ll \log_2 W)\,|\,j \;\in\; \bits^{d}, \qquad
  d = t + \log_2 W ,
\]
with $b$ the row, $c$ the column, $j$ the bit position inside the cell.
The $t$ row variables are folded with integer weights, the $s$ column
variables are read off in the clear.  We write $n = t + s$ for the
benchmark-scale parameter (the deployed shapes have $W = 1$, so $2^n$ is
the committed bit count) and $\K = \mathrm{GF}(2^{128})$ for the binding
field, with $\alpha \in \Kx$ a generator.

\paragraph{The claim the forest certifies.}
An integer fold cannot be read through the $\Ftwo$ commitment (parities
destroy magnitudes), so F2Z carries it in the exponent of $\K$: for one
chunk $l$ of the mod-$q$ opening, with public $\cw$-bit chunk weights
$W^{(l)}_b$ (limbs of the $\Fq$ row-$\eq$ weights,
$\cw = 127 - t - W$), the prover sends the per-column integer folds
\[
  u^{(l)}_c \;=\; \sum_b W^{(l)}_b \cdot \INT\!\big(D(b,c)\big)
  \;<\; 2^{\cw + t + W} = 2^{127},
\]
and must certify them against the commitment.  The certificate is a
grand product: each factor is \emph{bit-affine} in one committed bit,
\begin{equation}\label{eq:leaves}
  \Lambda^{(l)}[c][i] \;=\; 1 + M[c][i]\cdot\tau^{(l)}(i),
  \qquad
  \tau^{(l)}(i) \;=\; \alpha^{W^{(l)}_b 2^j} - 1
  \quad (i = (b \ll \log_2 W)|j),
\end{equation}
so that
$\prod_{i} \Lambda^{(l)}[c][i] = \alpha^{u^{(l)}_c}$.  (In $\K$,
$-1 = +1$; we keep the $-1$ form for readability.)  The forest's output
is a reduction of all $2^s$ product claims to \emph{one} evaluation
claim on the leaf MLE $\smash{\mle{\Lambda}}$, which
\S\ref{sec:discharge} converts into a $\K$-linear functional of the
committed bits and hands to the ring-switch + Ligerito opener.  Range
checks on the sent $u$'s and the generator property of $\alpha$ make the
exponent map injective on the values that occur, so the products bind
the integers; that chain of checks is the verifier note's
\S\S5--7, and we do not repeat it here.

\paragraph{Why performance of this phase is the whole game.}
Time is forest-dominated: the \code{forest+presum} phase is 81\,\% of
prove at $n = 16$, rising to 92--96\,\% at $n \ge 26$
(\S\ref{sec:costs}).  Proof \emph{bytes}, by contrast, are dominated by
the Ligerito opening --- the forest side stays at 8--51\,KiB through
$n = 28$.  So the forest is engineered for prover time and memory, and
its transcript is engineered to stay small; both axes are visible in the
design below.

\paragraph{Scope.}
Everything in this note is the code of this repository: the merged
forest (\code{src/merged\_forest.rs}), the eq-factored driver
(\code{src/piop/sumcheck/eq\_factored.rs}), and the table builders
(\code{src/pcs.rs}).  Function names are given so claims can be checked
against source; line numbers are avoided (the files drift).  The
Ligerito opener is out of scope (black box, as in the verifier note).

% ===========================================================================
\section{One forest, one transcript: the merged statement}
\label{sec:merged}

A naive realisation runs one GKR product tree per column and carries,
per tree per layer, a pair of child evaluations --- $2\cdot 2^s\cdot d$
$\K$-elements of proof, ${\approx}512$\,KiB at $n = 26$.  The merged
forest removes the $2^s$ factor from the payload by treating the tree
index as \emph{more MLE variables}.

\paragraph{The layer family.}
Index the values of tree $c$ at depth $\ell$ (the root is depth $0$,
leaves depth $d$) by $x \in \bits^\ell$, and define one multilinear
family per depth,
\[
  L_\ell(x, c), \qquad (x, c) \in \bits^\ell \times \bits^s,
\]
with $x$ occupying the \emph{low} index bits and $c$ the high ones.
The product recursion pairs children on the \emph{top} in-tree bit:
\begin{equation}\label{eq:recursion}
  L_\ell(x, c) \;=\; L_{\ell+1}(x, 0, c)\cdot L_{\ell+1}(x, 1, c),
\end{equation}
i.e.\ the two children of node $x$ sit at positions $x$ and
$x + 2^{\ell}$ of level $\ell + 1$.  This convention is load-bearing for
the lazy schedules: every level is stored and consumed as a top-bit
split pair of halves $(E, O)$ ---
$E_c(x) = L_{\ell+1}(x, 0, c)$, $O_c(x) = L_{\ell+1}(x, 1, c)$ --- and
the leaf halves are exactly ``first $2^{d-1}$ bits, last $2^{d-1}$
bits'' of a column, which is what lets pair products be table lookups
(\S\ref{sec:lazy}).  $L_0(\cdot, c) = R_c$ are the roots and
$L_d = \Lambda$ the leaves.

\paragraph{Entry: the roots are derived, absorbed, and evaluated ---
never sent per-tree.}
The prover's product-tree build produces the roots $R_c$ as a byproduct;
the proof object carries no root field.  On the verifier's side the
roots are \emph{recomputed} as $R_c = \alpha^{u_c}$ from the sent folds
(fixed-base comb, \code{FixedBasePow}) --- for an honest prover these
are the same values, and that equality is precisely the claim being
proven.  Both sides then perform the same two transcript actions
(\code{drive\_grouped} / \code{verify\_merged\_forest}):
\[
  \textbf{absorb}(\mathtt{0x30}: R_0, \dots, R_{2^s-1}),
  \qquad
  \zeta \leftarrow \textbf{squeeze}(\K^s),
\]
and each computes the entry claim itself,
$C_0 = \mle{R}(\zeta)$ (\code{mle\_at}, $O(2^s)$ work).  The absorb is
the binding: a prover whose subsequent messages were generated against
different roots faces challenges derived from \emph{these} roots.  Note
the prover-side integer folds $u_c$ (\code{fold\_values\_bits}, a
popcount-style pass over the packed rows) are computed independently of
the forest and enter the proof stream as the $\code{us}$ vector; the
forest itself never touches them.

\paragraph{Per-layer payload.}
Each of the $d$ layers contributes one \code{MergedLayer} to the proof:
an optional phase-A sumcheck ($\ell$ rounds, absent at $\ell = 0$), a
phase-B sumcheck ($s$ rounds), and a closing child pair --- $O(\ell+s)$
$\K$-elements.  Summed over layers this is
${\approx}\,48\big(\tfrac{d(d-1)}{2} + ds\big)$ bytes of round messages
plus $32d$ bytes of pairs (degree-3 rounds carry three 16-byte tail
values each, \S\ref{sec:driver}); at the deployed shapes the
\emph{sent folds} $16\cdot 2^s$\,B are the larger forest-side item
(e.g.\ $n = 28$, $s = 11$: 32\,KiB of folds vs.\ ${\approx}16$\,KiB of
forest messages, 50.7\,KiB forest-side total).

% ===========================================================================
\section{The layer protocol}
\label{sec:layers}

The verifier maintains a running claim $C_\ell = \mle{L_\ell}(z_x, z_c)$
starting from $(z_x, z_c) = (\varepsilon, \zeta)$.  By
\eqref{eq:recursion} and multilinearity, the layer-$\ell$ claim factors
coordinate-wise:
\begin{align}
  C_\ell
  \;&=\; \sum_{x \in \bits^\ell}\sum_{c \in \bits^s}
        \eq(x, z_x)\,\eq(c, z_c)\;
        \underbrace{L_{\ell+1}(x,0,c)}_{E_c(x)}\,
        \underbrace{L_{\ell+1}(x,1,c)}_{O_c(x)} \notag\\
  \;&=\; \sum_{c} \eq(c, z_c)
        \Big[\sum_{x} \eq(x, z_x)\, E_c(x)\, O_c(x)\Big].
        \label{eq:layer}
\end{align}
The bracketed inner sums share the point $z_x$ and differ only in the
tree; that shape is exactly one \emph{shared-$q$ multi-group} instance
of the eq-factored driver (\S\ref{sec:driver}), and the layer is
discharged in two phases plus a line step
(\code{drive\_grouped}):

\begin{enumerate}[leftmargin=2.2em, itemsep=2pt]
\item \textbf{Phase A} (in-tree variables; $\ell \ge 1$ only).  The
  $2^s$ trees enter as groups: group $c$ has the shared eq point
  $q = z_x$, the per-group scale $\eq(c, z_c)$ (a precomputed
  $O(2^s)$ table re-built per layer), and the single buffer pair
  $(E_c, O_c)$.  One driver call
  (\code{prove\_eq\_inner\_sumcheck\_mixed\_pre}) runs $\ell$ rounds of
  a degree-3 sumcheck binding $x \to r_x \in \K^\ell$, and returns the
  per-tree finals
  $\big(\mle{E_c}(r_x), \mle{O_c}(r_x)\big)$.  These finals are
  \emph{prover-internal}: nothing per-tree is absorbed at any point.
\item \textbf{Phase B} (tree-index variables).  The residual
  $\sum_c \eq(c, z_c)\,\mle{E_c}(r_x)\,\mle{O_c}(r_x)$ is itself an
  eq-weighted inner product over $c$ --- the per-tree finals, read as
  two length-$2^s$ multilinear tables $E(\cdot), O(\cdot)$, feed one
  more (tiny, dense) driver call with $q = z_c$, binding
  $c \to r_c \in \K^s$ and closing on the child pair
  \[
    (p, q) \;=\;
    \big(\mle{L_{\ell+1}}(r_x, 0, r_c),\;
         \mle{L_{\ell+1}}(r_x, 1, r_c)\big).
  \]
\item \textbf{Line step.}  $\textbf{absorb}(\mathtt{0x32}: p, q)$,
  $\lambda_\ell \leftarrow \textbf{squeeze}(\K)$, and the claim advances
  along the line through the two children:
  \[
    C_{\ell+1} \;=\; p + \lambda_\ell\,(p + q), \qquad
    z_x \leftarrow r_x\,\|\,\lambda_\ell, \qquad
    z_c \leftarrow r_c .
  \]
  (The code calls $\lambda_\ell$ \code{mu}; renamed to keep $\mu$ for
  the discharge residual, as in the verifier note.)
\end{enumerate}

Layer $0$ degenerates: there are no in-tree variables yet, phase A is
skipped (\code{sc\_x = None}), and the pair vectors fed to phase B are
simply each root's two children --- level $1$, two values per tree.
After $d$ layers the \emph{exit claim} is
\[
  e_d \;=\; C_d \;\overset{\text{honest}}{=}\;
  \mle{\Lambda}(z_x, z_c),
  \qquad z = (z_x, z_c) \in \K^{d+s},
\]
pinned by the \code{merged\_forest\_roundtrips} test against
\code{mle\_at(leaves, $z$)}.  The verifier-side chaining checks (phase-A
claimed sum $= C_\ell$, the factored-$\eq$ linkage
$E_A = \eq(r_x, z_x)\cdot\code{sc\_c.claimed\_sum}$, the closing check
$E_B = \eq(r_c, z_c)\cdot p\cdot q$) are [V5]--[V11] of the verifier
note and mirror this structure exactly.

\begin{remark}[Why two phases]
A single $(\ell + s)$-variable sumcheck over the merged index would be
sound too, but would need the $\eq(c, z_c)$ weights woven into every
round over a buffer of size $2^{\ell + s}$.  The split lets phase A
treat $\eq(c, z_c)$ as a per-group \emph{scalar} (it seeds the group's
prefix scalar, \S\ref{sec:driver}, at zero marginal cost), keeps the
per-tree buffers independent --- which is what the lazy bit-driven
rounds and the across-tree parallelism need --- and makes phase B
negligible ($2^s$-sized buffers).  The transcript cost is one extra
sumcheck header per layer.
\end{remark}

% ===========================================================================
\section{The eq-factored driver}
\label{sec:driver}

Every phase of every layer runs on one prover,
\code{prove\_eq\_inner\_sumcheck\_mixed\_pre}
(\code{src/piop/sumcheck/eq\_factored.rs}), which proves sums of the
shape
\[
  \sum_{x} \sum_{g} \eq(x; q_g)\cdot \mathrm{scale}_g \cdot
  L_g(x)\,R_g(x)
\]
--- an eq-weighted sum of inner products of multilinear pairs --- and is
\textbf{byte-identical} to the generic \code{MLSumcheck} prover run over
the materialised comb $[\eq_1, \dots, L\text{'s}, R\text{'s}]$ with the
degree-3 combiner $\sum_g \eq_g\, L_g R_g$: same round polynomials, same
transcript operations (header, tail absorbs, post-draw challenge
re-absorb), same proof layout.  The generic verifier verifies it
unchanged --- there is no bespoke verification path to audit.

\paragraph{Never materialise $\eq$.}
Round $j$ (binding the \emph{low} variable, as the generic prover does)
writes each group's round message as
\[
  P_j(c) \;=\; \sum_g A_{g,j}\;\eq_1\!\big(c;\, q_g[j-1]\big)\;
  H_{g,j}(c),
  \qquad
  H_{g,j}(c) \;=\; \sum_b V_{g,j}[b]\; L_g(c, b)\, R_g(c, b),
\]
where $L_g(c, b) = L_g[2b] + c\,(L_g[2b{+}1] - L_g[2b])$ is the
fold-in-progress buffer at suffix slot $b$, and the $\eq$ weight has
been split into three parts, none of them a materialised multiplicand:
\begin{itemize}[itemsep=1pt]
\item the \emph{prefix scalar}
  $A_{g,j} = \mathrm{scale}_g\cdot\prod_{i<j-1}\eq_1(\rho_{i+1};
  q_g[i])$, one running $\K$-element per group, updated per round ---
  this is where phase A's $\eq(c, z_c)$ group scale rides for free;
\item the current-variable factor $\eq_1(c; q_g[j-1])$, degree 1 in the
  round variable;
\item the \emph{suffix tensor}
  $V_{g,j}[b] = \prod_{i\ge j}\eq_1(b_{i-j}; q_g[i])$, precomputed
  back-to-front for all rounds in $O(2^k)$ total (\code{suffix\_tensors})
  --- the cost of one $\eq$ build, and built \emph{once} for the whole
  phase-A call, since all trees share $q = z_x$ (asserted:
  bit-selected groups require the shared point).
\end{itemize}
$H$ is quadratic and the $\eq_1$ factor linear, so $P_j$ has degree 3,
matching the pinned verifier degree.

\paragraph{Coefficient-space accumulation, wide sums, free fourth node.}
The per-slot inner loop accumulates the \emph{coefficients}
$(A_0, A_1, A_2)$ of $H$ in the round variable rather than its node
values, so the Lagrange-node multiplies stay out of the hot loop; the
single-pair case costs 5 field multiplies per slot (the suffix weight is
folded into $L$: $w{\cdot}l_0$, $w{\cdot}l_1$, then three products
$\smash{(w l_0)r_0}$, $\smash{(w l_1)r_1}$,
$\Delta L\cdot\Delta R$).  The three products feed
\code{WideMulAcc} accumulators: unreduced 256-bit carry-less sums with
\emph{one} polynomial reduction per accumulator per round --- exact,
because the reduction is $\Ftwo$-linear.  Round messages are sent as
tail evaluations $P_j(1), P_j(2), P_j(3)$ only; the verifier
reconstructs $P_j(0)$ from the chained sum (the claimed-sum trick
\cite{Gruen24}).  The fourth node is free in characteristic 2: the
evaluation points $\{0, 1, X, X{+}1\}$ (bit-pattern convention) form an
$\Ftwo$-affine 2-flat, and over such a flat every polynomial of degree
$\le 2$ sums to zero, so $H(X{+}1) = H(0) + H(1) + H(X)$.  The
detection is exact at runtime ($1 + 1 = 0$), keeping the driver
field-generic.  On aarch64 the whole-buffer message and fold passes
dispatch to fused NEON kernels (\code{eqf\_single\_pair\_round} and
friends; \code{F2Z\_EQF\_NOKERNEL=1} forces the generic loops for
A/B diagnosis).

\paragraph{Pass fusion.}
By default the driver defers each round's fold and runs it \emph{fused
into the next round's message pass}: one read of the unfolded buffers
replaces a fold-write pass plus a message-read pass.  Measured
${\approx}10\,\%$ of prove at $n = 26$--$28$; byte-identical (same field
values in the same transcript order --- only the physical pass structure
changes); \code{F2Z\_EQF\_FUSE=0} restores the eager two-pass path.
The same idea appears one level up in the forest's JIT layers: the
buffer-\emph{generation} pass accumulates the round-1 coefficients as it
writes (\code{dense\_jit\_fused\_round1}, forwarded through the driver's
\code{pre\_round1} hook), so a JIT-generated level is written once and
first \emph{read} by round 2's fused fold+message pass
(\code{F2Z\_JIT\_R1=0} opts out).  XOR accumulation is order-free and
reduction $\Ftwo$-linear, hence generation-order independence of the
reduced coefficients --- the fusion is again an exact identity.

% ===========================================================================
\section{Leaves from bits: the lazy schedules}
\label{sec:lazy}

The eager reference prover (\code{prove\_merged\_forest}) materialises
the $2^{d+s}$ leaves and every product level above them.  The lazy
prover (\code{prove\_merged\_forest\_lazy}) produces a byte-identical
transcript while never materialising the leaf layer or the bottom
product levels at all.  Everything rests on the observation that at and
near the leaves, the forest's values have so little entropy that they
can be \emph{gathered from small shared tables indexed by committed
bits} instead of computed and stored.

\subsection{What the prover holds}
\label{sec:tables}

\paragraph{Bits.}
The commitment hint stores the bit matrix column-lane--packed (64
columns per \code{u64} word).  One transposition pass
(\code{extract\_column\_bit\_halves}, blocks of $64{\times}64$
bit-transposes) yields, per tree $c$, the top-bit-split halves of its
leaf bits:
\[
  \code{lbits}_c = M[c][0..2^{d-1}], \qquad
  \code{rbits}_c = M[c][2^{d-1}..2^d],
\]
packed 64 per word --- a few KiB per tree.  These two bit arrays drive
\emph{every} lazy layer of tree $c$'s bottom.

\paragraph{$\alpha$-power chains.}
Per chunk, \code{chunk\_pow2\_table} builds
$v(i) = \alpha^{W^{(l)}_b 2^j}$ for all $i = (b \ll \log_2 W)|j$: one
fixed-base comb for the shared $\alpha$-squaring chain
(\code{FixedBasePow}, ${\approx}8\times$ fewer multiplies than
square-and-multiply per exponent), then a per-row squaring chain for the
$2^j$ place values.

\paragraph{Shared value tables.}
From the chains, four table families are built once per weight set
(``$\tau$ set'') and shared by all $2^s$ trees:
\begin{itemize}[itemsep=1pt]
\item \emph{leaf $\tau$ halves} (\code{leaf\_tau\_halves}):
  $\tau(i) = v(i) + 1$, split like the bits --- the affine coefficients
  of \eqref{eq:leaves};
\item the \emph{pair table} (\code{layer1\_pair\_table}):
  $\code{pair}[i] = v(i)\cdot v(i + 2^{d-1})$, the only case of a
  level-$(d{-}1)$ value that is not already a table entry;
\item the \emph{4-case tables} $\code{te}, \code{to}$: position $y$ of
  level $d{-}1$ pairs leaves $(y,\; y + 2^{d-1})$, so its value is one
  of $\{1,\, v(y),\, v(y{+}2^{d-1}),\, \code{pair}[y]\}$, selected by
  two bits.  \code{te} tabulates the four cases for the E-half positions
  $y < 2^{d-2}$, \code{to} for the O-half positions (offset baked in);
\item the \emph{16-case table} $T4$:
  $T4[(y \ll 4)|(c_E \ll 2)|c_O] = \code{te}[4y{+}c_E]\cdot
  \code{to}[4y{+}c_O]$ --- all 16 possible values of the
  \emph{4-leaf product} at level $d{-}2$, position $y$, selected by four
  committed bits.
\end{itemize}
Costs are $O(2^d)$ per set --- $2^s$-fold amortised --- and the
$T4$ build uses an indexed collect with an in-place flatten (a parallel
nested \code{flatten} here measured ${\approx}836\times$ slower than the
serial arithmetic; upstream fix \code{e19b0e1}).

\subsection{The level ladder}
\label{sec:ladder}

Recall layer $\ell$'s phase A consumes level $\ell{+}1$.  The lazy
prover sources each level in one of three ways: from the \emph{stored
chain} (built bottom-up once, level-major, by \code{build\_levels}),
\emph{JIT} (regenerated from bits and tables while its layer runs, then
dropped), or \emph{implicitly} (never as $\K$-values: the layer's first
round(s) run case-LUT kernels straight off the bits,
\S\ref{sec:lutrounds}).  Two schedules exist, named for the size of the
stored chain relative to the $2^d$-leaf layer; both are pinned
byte-identical to the eager forest.

\paragraph{The L/4 schedule (default; depth $\ge 4$).}
\begin{center}
\small
\begin{tabular}{llll}
\toprule
layer $\ell$ & consumes level & sourced & phase-A entry \\
\midrule
$0 \dots d{-}4$ & $1 \dots d{-}3$ & stored chain & \code{Dense} \\
$d{-}3$ & $d{-}2$ & JIT: one $T4$ gather/value & \code{Dense} (round 1 fused) \\
$d{-}2$ & $d{-}1$ & implicit & \code{Pair3Bits} over $(\code{te},\code{to})$ \\
$d{-}1$ & $d$ (leaves) & implicit & \code{Leaf3Bits} over $\tau$ \\
\bottomrule
\end{tabular}
\end{center}
The stored chain is built per tree directly at level $d{-}3$ --- each
value the product of two paired $T4$ gathers --- so levels $d{-}2$,
$d{-}1$, $d$ never exist at build time, and the chain then folds upward
($\sum_{\ell \le d-3} 2^\ell \approx 2^{d-2}$ values per tree, a quarter
of the leaf count; hence the name).  Reading the bits word-wise during
these builds costs 4 \code{u64} loads per 64 values,
${\approx}64\times$ less read traffic than per-bit gathers over the
column-lane store.  At depth 4 exactly (and for every depth under the
\code{F2Z\_LUT3=0} opt-out) the bottom two layers run the shallower
\code{Pair2Bits} / \code{Leaf2Bits} variants; depth 3 keeps only the
leaf layer bit-driven (\code{LeafBits}) with the chain built from
level-1 halves by a zero-multiply case copy
(\code{build\_column\_layer1\_halves}); depth $< 3$ falls back to the
eager prover (the LUT rounds need $k \ge 2$).

\paragraph{The L/8 schedule (\code{F2\_FOREST\_SCHEDULE=l8}; depth $\ge 5$).}
\begin{center}
\small
\begin{tabular}{llll}
\toprule
layer $\ell$ & consumes level & sourced & phase-A entry \\
\midrule
$0 \dots d{-}5$ & $1 \dots d{-}4$ & stored chain & \code{Dense} \\
$d{-}4$ & $d{-}3$ & JIT: two $T4$ gathers + 1 mul /value & \code{Dense} (round 1 fused) \\
$d{-}3$ & $d{-}2$ & implicit & \code{T4Bits} over $T4$ \\
$d{-}2$ & $d{-}1$ & implicit & \code{Pair3Bits} over $(\code{te},\code{to})$ \\
$d{-}1$ & $d$ & implicit & \code{Leaf3Bits} over $\tau$ \\
\bottomrule
\end{tabular}
\end{center}
Every stage's resident set is ${\approx}2^{d-3}$ values per tree ---
an eighth of the leaf count --- at a measured $+5$--$13\,\%$ prove cost
at the 2-column model shapes; the schedule exists for memory-bound runs
and is \emph{required} at $n \ge 31$ on a 16\,GB box
(\S\ref{sec:costs}).

\subsection{Inside the case-LUT rounds}
\label{sec:lutrounds}

The bit-driven layers are where ``lazy'' earns its keep: their input
levels are never materialised because the first one, two, or three
sumcheck rounds of the layer can be computed \emph{directly} from
$(\code{lbits}, \code{rbits})$ and per-position case tables.  All of the
following are exact char-2 identities --- the transcripts they produce
are bit-for-bit those of dense buffers.

\paragraph{The leaf round (\code{LeafBits} family).}
For the leaf layer, entry $i$ of the phase-A pair is \emph{defined} as
$L[i] = 1 + \code{lbits}[i]\cdot\tau_L[i]$ (and $R$ from
$\code{rbits}, \tau_R$).  Expanding round 1's per-slot contribution
$w_b\,L(c)R(c)$ over the four selecting bits, the $1$'s cancel
(char 2) into subset sums of eight tree-shared $w\tau$-product tables
(\code{build\_leaf\_tables}: the 4-case singles/pair tables
$t_{a_0}, t_{a_1}$ and a $\Delta L\,\Delta R$ table), so a
\emph{group's} round-1 message costs \textbf{zero multiplications}: per
slot, one unconditional indexed load and add per accumulator (3 loads
$+$ 3 adds), branchless on the bit nibble.  The tables themselves cost
8 multiplies $+$ ${\approx}20$ adds per slot to build --- once,
amortised over all $2^s$ trees.  The $\Delta\Delta$ table has two
interchangeable forms, chosen by footprint
(\code{F2Z\_LEAF\_A2\_FACTORED}, auto at $2^{k-1} \ge 2^{15}$): a
16-case precombined table (one load per slot; wins while L2-resident)
or the four raw cross products with branchless masked adds
($4\times$ smaller; at $n = 30$ the leaf round drops $2.1\times$,
$307 \to 147$\,ms, prove $-6.5\,\%$).

\paragraph{Folding without materialising.}
Round 1's fold must produce round 2's buffers.  For a bit-affine leaf
side this is again a 4-case table: the folded entry is
$1 + m_0(1{+}\rho_1)\tau_0 + m_1\rho_1\tau_1$, precombined per position
over the adjacent bit pair $(m_0, m_1)$
(\code{build\_leaf\_fold\_tables}), so the fold is one indexed load per
entry.  The deeper variants iterate the idea instead of executing it:
\begin{itemize}[itemsep=1pt]
\item \code{Leaf2Bits} ($k \ge 3$): round 1's fold \emph{keeps the
  bits} and stashes its fold tables as round 2's \emph{value} tables ---
  which are structurally a $(\code{te}, \code{to})$ pair, so round 2
  runs the 16-case product-layer round below, and only round 2's fold
  materialises dense buffers ($2^{k-2}$ per side);
\item \code{Leaf3Bits} ($k \ge 4$, the default): one round deeper
  still; round 3 reads the stashed 16-case fold tables \emph{inline}
  (one aligned byte load per side per slot --- the measured lesson being
  that multiplies are free and only skipped bytes pay), and dense
  buffers appear only at round 3's fold ($2^{k-3}$ per side).
\end{itemize}

\paragraph{The product-layer round (\code{Pair2Bits}/\code{Pair3Bits}).}
One level up, entry $y$ of the pair is the 2-bit select
$\code{te}[4y + (m_{\mathrm{lo}} | m_{\mathrm{hi}}{\ll}1)]$ (and
\code{to} at the offset position), over the \emph{same} leaf-bit halves.
The default round body is the \emph{factored} form
(\code{F2Z\_PAIR2\_FACTORED}): only the suffix-weighted table
$\code{wte}[i] = V_1[i \gg 3]\cdot\code{te}[i]$ is built, and each slot
does three wide multiplies against \code{wte}/\code{to} entries ---
trading one multiply for a $4\times$ smaller, L2-resident table set
(the precombined 16-case alternative carries $64\cdot2^k$ entries,
16\,MiB at deployed shapes, and its four gather streams' misses
dominate the round).  \code{Pair3Bits} again stashes round 1's 16-case
fold tables as round 2's inline value tables and materialises at round
2's fold.

\paragraph{The 4-leaf round (\code{T4Bits}, L/8 only).}
One further level up, entry $y$ is a single 16-case select from the
shared $T4$ --- one gather per value, values read inline in round 1,
dense buffers at its fold.  The layer's input level is never stored
\emph{nor} regenerated.

\paragraph{Prefetch.}
The LUT gathers are data-dependent (committed bits pick the line), which
defeats the hardware prefetcher, but the indices are cheaply
recomputable ahead: both the stash-gather rounds
(\code{F2Z\_LUT\_PRFM}) and the $T4$ build/JIT sites
(\code{F2Z\_T4\_PRFM}) issue \code{prfm pldl1keep} 16 slots ahead,
size-gated to the DRAM-scale shapes (they \emph{lose} at L2-resident
sizes; measured $-4.8\,\%$ and $-5.3\,\%$ prove at $n = 30$
respectively).  Semantically inert.

\subsection{Why byte-identity holds, and how it is pinned}
\label{sec:identity}

Every lazy device above is a re-\emph{association} of the same field
arithmetic: case tables precombine subexpressions that the dense path
would compute per entry (field addition is exact and commutative);
char-2 cancellations ($1 + 1 = 0$, $v_1 - v_0$'s ones cancelling in
folds) are identities; wide-vs-narrow accumulation reduces to identical
values because the reduction is $\Ftwo$-linear; fusion only reorders
physical passes.  Consequently the lazy prover does not merely
``verify the same'' --- it emits the \emph{same transcript}.
\code{lazy\_matches\_eager} pins roots, every layer's sumcheck messages,
pairs, exit claim, and the post-run transcript state against the eager
prover across depths 3--9 (covering every fallback edge and both the
scalar and word-wise $T4$ paths), $W \in \{1, 4\}$, and \emph{both}
schedules; \code{lazy\_multi\_matches\_eager} does the same for the
batched forest (\S\ref{sec:variants}).  This is why the knob inventory
of Table~\ref{tab:knobs} is safe: every knob changes scheduling, not
semantics.

\begin{table}[ht]
\centering\small
\begin{tabular}{lll}
\toprule
Knob & Default & Effect (always byte-identical) \\
\midrule
\code{F2\_FOREST\_SCHEDULE=l8} & off (L/4) & L/8 memory schedule (\S\ref{sec:ladder}) \\
\code{F2Z\_LUT3=0} & on & shallower \code{Pair2}/\code{Leaf2} LUT prefixes \\
\code{F2Z\_EQF\_FUSE=0} & on & eager two-pass fold (no pass fusion) \\
\code{F2Z\_JIT\_R1=0} & on & unfused JIT round-1 message pass \\
\code{F2Z\_LEAF\_A2\_FACTORED=0/1} & auto ($\ge 2^{15}$) & $\Delta\Delta$ table form, leaf round \\
\code{F2Z\_PAIR2\_FACTORED=0} & on & precombined 16-case pair tables \\
\code{F2Z\_LUT\_PRFM=0/1} & auto ($\ge 2^{14}$) & stash-gather software prefetch \\
\code{F2Z\_T4\_PRFM=0/1} & auto ($\ge 16$\,MiB) & $T4$ build/JIT software prefetch \\
\code{F2Z\_EQF\_NOKERNEL=1} & off & bypass fused NEON kernels (diagnostic) \\
\bottomrule
\end{tabular}
\caption{The forest's scheduling knobs.  Each is read once per
prove/process; none changes the proof bytes.}
\label{tab:knobs}
\end{table}

% ===========================================================================
\section{Memory and cost}
\label{sec:costs}

\paragraph{Memory model.}
With 16-byte $\K$-elements and $W = 1$ (so $2^{d+s} = 2^n$ leaves), the
resident forest set is capped by the largest stage of the ladder:
\[
  \text{L/4:}\quad
  {\approx}\,2^{n-2}\cdot 16\,\mathrm{B} = 2^n\cdot 4\,\mathrm{B},
  \qquad
  \text{L/8:}\quad
  {\approx}\,2^{n-3}\cdot 16\,\mathrm{B} = 2^n\cdot 2\,\mathrm{B},
\]
(the stored chain and the JIT level are the same size and not
simultaneously live: the chain is consumed front-first,
\code{mem::take} per layer, before the JIT layer runs).  On top sit the
commit hint (${\approx}\,2^n\cdot 0.7$\,B: packed rows/columns,
codeword, Merkle tree), the transposed bit halves ($2^n/8$\,B), and the
shared tables ($O(2^d)$ per set; the largest, $T4$, is
$16\cdot 2^{d-2}$ entries --- 8\,MiB at $d = 17$).  The model matches
the measured peaks below to within ${\approx}5\,\%$; it is why the
16\,GB ceiling sits at $n = 32$ (under \code{l8}) and why $n \ge 31$
requires \code{l8} at all.

\paragraph{Measured (Apple M4 16\,GB, LTO + \code{unchecked}, fusion
defaults, rate-1/2 $k{=}4$ configs, medians; 2026-07-18).}
The \code{forest+presum} column is the phase this note describes (the
$\alpha$-power tables, the forest proper, and the pre-sumcheck of
\S\ref{sec:discharge}); its share of prove rises from 81\,\% to
${\approx}96\,\%$.
\begin{center}
\small
\begin{tabular}{rlrrrrrl}
\toprule
$n$ & $(t, s)$ & forest+presum & lig.\ open & prove & verify & peak & sched \\
\midrule
16 & (10, 6)  & 3.25\,ms & 0.72\,ms & 4.03\,ms & 1.21\,ms & 0.69\,MB & L/4 \\
20 & (13, 7)  & 9.13\,ms & 1.62\,ms & 10.9\,ms & 1.83\,ms & 7.7\,MB  & L/4 \\
24 & (15, 9)  & 42.5\,ms & 21.5\,ms & 62.1\,ms & 2.48\,ms & 85.0\,MB & L/4 \\
26 & (16, 10) & 131\,ms  & 25.3\,ms & 159\,ms  & 3.14\,ms & 325\,MB  & L/4 \\
28 & (17, 11) & 530\,ms  & 53.6\,ms & 579\,ms  & 4.41\,ms & 1.27\,GB & L/4 \\
30$^\dagger$ & (18, 12) & 2.98\,s & 139\,ms & 3.15\,s & 5.55\,ms & 4.95\,GB & L/4 \\
32$^\dagger$ & (19, 13) & ${\sim}96\,\%$ & ${\sim}4\,\%$ & 29.8\,s & 11.8\,ms & 11.5\,GB & \code{l8} \\
\bottomrule
\end{tabular}
\end{center}
($^\dagger$: the older $k{=}6$ generation, pending a memory-fresh
re-measure; full tables and the measurement protocol in
\code{README.md}.)  Proof bytes stay forest-light: forest-side
7.7\,KiB ($n{=}16$) $\to$ 50.7\,KiB ($n{=}28$) $\to$ 151.8\,KiB
($n{=}32$), of which the sent folds are $16\cdot2^s$\,B
(\S\ref{sec:merged}); the Ligerito blob dominates the stream
throughout.  Threading: the build is parallel across trees, round
bodies across groups (with a minimum batch so tiny late rounds amortise
dispatch), phase B and the upper layers effectively serial; the
MT/ST prove ratio grows from ${\approx}1.6\times$ at $n = 22$ to
${\approx}3.6\times$ at $n = 28$.  All deployed shapes run a
\emph{single} chunk (with $q = 2^{100}{-}15$, $\cw = 126 - t \ge 100$
for $t \le 26$); chunking multiplies the forest count only at
$t + W > 27$.

% ===========================================================================
\section{The exit claim and its discharge}
\label{sec:discharge}

The forest ends holding $e_d \approx \mle{\Lambda}(z_x, z_c)$ --- still
entangled with the $\alpha$-power weights.  Two further steps
(\code{prove\_int\_eval\_merged\_common}) turn it into something the
opener understands; they are the ``presum'' half of the
\code{forest+presum} phase.

\paragraph{From the exit claim to a bit-MLE claim.}
By \eqref{eq:leaves} and multilinearity (the all-ones table extends to
$1$):
\[
  e_d - 1
  \;=\; \sum_{i, c} \eq(i, z_x)\,\eq(c, z_c)\, M[c][i]\,\tau(i)
  \;=\; \sum_{i \in \bits^d} q_{\mathrm{rowbit}}[i]\cdot m[i],
\]
with $q_{\mathrm{rowbit}}[i] = \eq(i, z_x)\,\tau(i)$
(\code{row\_bit\_weights} at $\rho = z_x$; its $2^t$
${\approx}q_{\mathrm{bits}}$-bit exponentiations go through the same
fixed-base comb) and $m[i] = \sum_c \eq(c, z_c)\,M[c][i]$ the
$z_c$-combined row (\code{xi\_combined\_rows}).  A $d$-round
\emph{degree-2} sumcheck over this product
(\code{MultiDegreeSumcheck}, one group) binds $i \to r^*$; dividing the
final expected evaluation by the public
$\smash{\mle{q_{\mathrm{rowbit}}}}(r^*)$ leaves
$\mu = \mle{M}(r^*, z_c)$ --- a pure bit-MLE evaluation at
$(r^* \| z_c)$, every $\alpha$-dependent factor stripped to the
verifier's side.  From here the ring switch sends the 128 partial
evaluations $s_v$, the $L$ chunks' residuals are $\eta$-combined, and
one recursive Ligerito call certifies the batch against the commitment
--- verifier note \S6, out of scope here.

\paragraph{The chunked flow.}
\code{prove\_mle\_eval\_mod\_q\_ligerito} runs the full
forest~$+$~presum \emph{per chunk}, sequentially on one transcript
(chunk $l$'s $\alpha$-power tables are built and dropped inside its
run), then does all ring-switch messages, then the single Ligerito
call.  Per chunk the proof stream carries: the \code{MergedForestProof},
the folds $u^{(l)}$, the pre-sumcheck, and $s_v$
(\code{proof\_codec.rs}).

% ===========================================================================
\section{Variants: the batched forest and the RLC families}
\label{sec:variants}

\paragraph{Multi-claim batching (\code{prove\_merged\_forest\_lazy\_multi}).}
$N$ same-shape claims --- each $2^s$ trees over its own bits and
$\tau$ chains --- run as \emph{one} forest of $N\cdot 2^s$ trees, tree
index $(\nu \ll s)\,|\,c$ with the claim index high, sharing every
layer's rounds and messages.  Per-claim table selection rides the
driver's \code{tau\_set} index (the kernels are untouched); table sets
are deduplicated by $\tau$-chain identity, so same-point claims share
one set.  $N$ must be a power of two; callers pad with zero-weight
dummies, whose $\tau$ chains are all $1$ and whose leaves are therefore
identically $1$ regardless of the bits they carry.  Both schedules
honoured; pinned byte-identical to the eager prover over the
concatenated leaf tables, dummy arm included.

\paragraph{RLC claim families (\textsc{experimental}).}
The mod-$q$ RLC families (\code{docs/rlc-family-note/}) replace the
bit-affine leaf with a $2^j$-case select on $j$ committed column bits:
$\Lambda[c][i] = \code{case\_pow}[i][m(i,c)]$, case $0 = \alpha^0 = 1$.
For $j = 2$ (\code{prove\_merged\_forest\_lazy\_rlc2}) this is a pure
\emph{rewiring} of the machinery above, one level shifted: the leaf
layer is structurally a product layer (\code{Pair3Bits} over
$\code{te}[4y{+}m] = \code{case\_pow}[y][m]$), level $d{-}1$ is a
\code{T4Bits} round over $\code{te}\cdot\code{to}$ products, level
$d{-}2$ is Dense JIT from paired gathers, and the chain tops at $d{-}3$
--- the L/4 shape with untouched kernels, hence byte-identical to the
eager forest over the same leaves (pinned).  For $j = 3, 4$
(\code{prove\_merged\_forest\_lazy\_rlc\_general}) no $2^j$-case leaf
kernels exist yet, so the bottom layers run Dense with JIT-generated
buffers gathered from a shared $2^{2j}$-case top-pair product table
$T2$; dedicated 8/16-case leaf-round kernels are the stated open lever.
(A Dense-JIT arm for the \emph{bit-affine} $j \ge 3$ forests was
measured slower than the eager path and closed.)

% ===========================================================================
\section{What to trust, and what this note is not}
\label{sec:trust}

The prover described here is exercised by: the both-schedule
byte-identity pins (\code{lazy\_matches\_eager},
\code{lazy\_multi\_matches\_eager} --- roots, every message, exit
claim, transcript state, across the depth-regime and $W$ matrix), the
GKR roundtrip and tamper rejections (\code{merged\_forest\_roundtrips}),
the mod-$q$ end-to-end roundtrips with range/generator/tamper arms, and
the NEON/scalar field-pipeline equivalence
(\code{neon\_mul\_matches\_scalar\_pipeline}).  Soundness of what the
forest \emph{binds} --- the chain from the absorbed roots through the
layer sumchecks to the exit claim, the exponent injectivity, and the
oracle handoff --- is argued in the verifier note and
\code{docs/DESIGN.md}; this note's claim is narrower and mechanical:
the lazy, merged, fused prover computes \emph{the same protocol} as the
eager reference, at $O(\ell + s)$ payload per layer,
${\approx}\,2^n\cdot4$\,B (or $2^n\cdot2$\,B) of working set, and
81--96\,\% of an F2Z prove.  The standard caveats carry over unchanged:
flock's Merkle domain separation is a flagged upstream item, the
ad-hoc Ligerito config regime ($m < 22$) is unaudited, and the scheme
is not zero-knowledge.

% ===========================================================================
\begin{thebibliography}{9}\small

\bibitem{GKR15} S.~Goldwasser, Y.~T.~Kalai, G.~N.~Rothblum.
\emph{Delegating computation: interactive proofs for muggles.}
J.~ACM 62(4), 2015 (STOC 2008).

\bibitem{Thaler13} J.~Thaler.
\emph{Time-optimal interactive proofs for circuit evaluation.}
CRYPTO 2013.  (The product-tree GKR specialisation the forest layers
instantiate.)

\bibitem{Libra} T.~Xie, J.~Zhang, Y.~Zhang, C.~Papamanthou, D.~Song.
\emph{Libra: succinct zero-knowledge proofs with optimal prover
computation.}  CRYPTO 2019.  (The eq-factored linear-time GKR prover
lineage.)

\bibitem{Gruen24} A.~Gruen.
\emph{Some improvements for the PIOP for ZeroCheck.}
IACR ePrint 2024/108.  (Tail-only round messages via the claimed-sum
reconstruction.)

\bibitem{Soukhanov} L.~Soukhanov.
\emph{The characteristic-2 field-switching note}, \S9.  (The
integer-MLE-over-$\Ftwo$ construction F2Z realises; internal.)

\end{thebibliography}

\end{document}
